\section{用益物权}
\begin{enumerate}
  \item 1 道单选题。
  \item 用他人之物让自己获益。支配的是他人之物的\zd{使用价值}。
  \item 建设用地使用权。
  \item \zd{土地承包经营权（农村）= 合同生效时 + 登记对抗 > 第三人}。流转时未经登记不得对抗善意第三人。
        \begin{enumerate}
          \item \zd{耕地}、草地、林地 \yjt 无偿取得，只发包给本村。村外：优先承包权；2/3 + 报乡镇政府批准。
          \item 四荒，荒山、荒沟、荒丘、荒滩 \yjt 有偿，招、拍、协。
        \end{enumerate}
  \item 宅基地使用权。
  \item \zd{居住权 = 书面有效合同/有效遗嘱 + 登记生效}。
        \begin{enumerate}
          \item 离婚经济帮助，\al{离婚不分家}；以房养老。
          \item 住宅上设置的用益物权：占有、使用、收益（相对禁止，有约从约）、处分（绝对禁止，转让 + 继承）。
          \item 可有偿，可无偿；\zd{有期限}。
        \end{enumerate}
  \item \zd{地役权}。供役地，需役地。合同生效时成立，登记对抗。
        \begin{enumerate}
          \item 产生方式：合同、意定、有偿、有期限。\al{海景房}。
          \item 从属性、不可分性。
          \item 和相邻关系区分。法定的，无期限。
        \end{enumerate}
\end{enumerate}
